\documentclass[a4paper, serif, 10pt]{moderncv}

%% moderncv
% style and color
\moderncvstyle{banking}
\moderncvcolor{black}

%% page layout/margins
\usepackage[top=2.5cm, bottom=2.5cm, left=1.5cm, right=1.5cm]{geometry}

\nopagenumbers{}

%% Babel config for Dutch language
% ref:
% - https://latex3.github.io/babel/guides/locale-dutch.html
\usepackage[dutch]{babel}

%% fontspec
\usepackage{fontspec}

\IfFontExistsTF{Linux Libertine}{
  \setmainfont[Ligatures={TeX, Common}, Numbers=OldStyle]{Linux Libertine}
}{}
\IfFontExistsTF{Linux Libertine O}{
  \setmainfont[Ligatures={TeX, Common}, Numbers=OldStyle]{Linux Libertine O}
}{}

%% CTeX
% CJK support, used to show CJK characters in the resume
%
% - fontset=none: disable builtin fontset but instead set the CJK font manually
% - heading=false: disable ctex heading
% - punct=kaiming: use kaiming punctuations styles for CJK
% - scheme=plain: use plain scheme, do not override \normalsize font size
% - space=auto: space settings for CJK characters
%
% ref:
% - http://ctan.mirrorcatalogs.com/language/chinese/ctex/ctex.pdf
\usepackage[UTF8, heading=false, punct=kaiming, scheme=plain, space=auto]{ctex}

\IfFontExistsTF{Noto Serif CJK SC}{
  \setCJKmainfont{Noto Serif CJK SC}
}{}
\IfFontExistsTF{Noto Sans CJK SC}{
  \setCJKsansfont{Noto Sans CJK SC}
}{}

%% line spacing
\usepackage{setspace}
\setstretch{1.125}

%% URL styling - use normal text instead of monospace
\urlstyle{same}

% Auto-underline all links
% Defer until \begin{document} so hyperref is already loaded
\AtBeginDocument{%
  \let\oldhref\href
  \renewcommand{\href}[2]{\oldhref{#1}{\underline{#2}}}%
}

%% Patch \httplink and \httpslink to handle full URLs
% The original moderncv macros always prepend http:// or https:// to URLs.
% This causes issues when the URL already has a protocol (e.g., https://example.com)
% resulting in malformed URLs like https://https://example.com.
% This patch checks if the URL already contains "://" and uses it directly if so.
% Note: We use \str_set:Nx (x-expansion) to expand macros like \@homepage first.
\ExplSyntaxOn
\str_new:N \g_yamlresume_url_str

\RenewDocumentCommand{\httpslink}{O{}m}{
  \str_set:Nx \g_yamlresume_url_str {#2}
  \str_if_in:NnTF \g_yamlresume_url_str {://}
    { \url{#2} }
    { \url{https://#2} }
}

\RenewDocumentCommand{\httplink}{O{}m}{
  \str_set:Nx \g_yamlresume_url_str {#2}
  \str_if_in:NnTF \g_yamlresume_url_str {://}
    { \url{#2} }
    { \url{http://#2} }
}
\ExplSyntaxOff

\name{Sanne van den Berg}{}
\title{Data-analist met passie voor datagedreven besluitvorming}
\phone[mobile]{+31 6 12345678}
\email{sanne.vandenberg@voorbeeld.nl}
\homepage{https://www.linkedin.com/in/sannevandenberg}

\address{Keizersgracht 123, Amsterdam, Noord-Holland, Nederland, 1016 CJ}{}{}

\extrainfo{{\small \faLinkedin}\ \href{https://www.linkedin.com/in/sannevandenberg}{@sannevandenberg} {} {} {} • {} {} {}
{\small \faGithub}\ \href{https://github.com/sannevandenberg}{@sannevandenberg}}

\begin{document}

\maketitle

\section{Basis Informatie}

\cvline{}{Data-analist met ruim 9 jaar ervaring in het omzetten van ruwe data naar
concrete inzichten voor financiële en commerciële teams. Ik combineer
sterke technische vaardigheden in SQL, Python en Power BI met een
gestructureerde, adviserende werkwijze.
%
\begin{itemize}
\item Ontwikkel en beheer dashboards en datapijplijnen voor diverse stakeholders
\item Ervaren in A/B-testen, statistische analyses en datamodellering
\item Sterke focus op datakwaliteit, reproduceerbaarheid en heldere rapportage
\item Werk graag op het snijvlak van data, business en technologie
\end{itemize}}
\section{Opleidingen}

\cventry{sep 2011–jul 2014}
        {Bachelor, Econometrie en Operationele Research, Score: 7.8}
        {Vrije Universiteit Amsterdam}
        {\url{https://www.vu.nl}}
        {}
        {\begin{itemize}
\item Afstudeerrichting toegepaste econometrie met nadruk op data-analyse en modellering
\item Scriptie over de voorspellende waarde van consumentenvertrouwen voor de detailhandel
\end{itemize}
\textbf{Cursussen}: Statistiek, Wiskundige economie, Operationele research, Econometrie}

\cventry{sep 2014–jul 2016}
        {Master, Data Science and Business Analytics, Score: 8.2}
        {Universiteit van Amsterdam}
        {\url{https://www.uva.nl}}
        {}
        {\begin{itemize}
\item Afgestudeerd met een 8,2 voor de scriptie over churnvoorspelling in de telecomsector
\item Specialisatie in voorspellende modellen en datagedreven besluitvorming
\end{itemize}
\textbf{Cursussen}: Machine Learning, Big Data Analytics, Datavisualisatie, Statistisch leren}
\section{Werkervaring}

\cventry{jan 2020–Heden}
        {Data-analist}
        {Finlytics B.V.}
        {\url{https://www.finlytics.nl}}
        {}
        {\begin{itemize}
\item Ontwikkel en beheer financiële dashboards in Power BI voor directie en businesscontrollers
\item Automatiseer data-extracties en transformaties met Python en dbt, waardoor rapportagetijd met 40\% afnam
\item Voer ad-hocanalyses uit en adviseer over kostenbesparende maatregelen
\item Begeleid junior analisten en stel richtlijnen op voor datakwaliteit en visualisaties
\end{itemize}
\textbf{Trefwoorden}: Power BI, SQL, Python, dbt, Financiële rapportages}

\cventry{okt 2017–dec 2019}
        {Data-analist}
        {RetailIQ}
        {\url{https://www.retailiq.nl}}
        {}
        {\begin{itemize}
\item Droeg bij aan de ontwikkeling van een klantsegmentatiemodel op basis van RFM-analyse
\item Analyseerde klantgedrag en conversiepaden in Google Analytics en BigQuery
\item Voerde A/B-tests uit op de webshop en rapporteerde over statistische significantie
\item Bouwde automatische datapijplijnen in Python voor dagelijkse verkooprapportages
\end{itemize}
\textbf{Trefwoorden}: BigQuery, Google Analytics, Klantsegmentatie, A/B-testen}

\cventry{sep 2016–sep 2017}
        {Junior data-analist}
        {Zorgdata Nederland}
        {\url{https://www.zorgdata.nl}}
        {}
        {\begin{itemize}
\item Verwerkte en controleerde zorgdatasets met SQL en SPSS
\item Maakte rapportages en grafieken voor interne en externe stakeholders
\item Ondersteunde senior onderzoekers bij datakwaliteitscontroles
\end{itemize}
\textbf{Trefwoorden}: SQL, SPSS, Datakwaliteit, Rapportage}
\section{Talen}

\cvline{Nederlands}{Moedertaal of tweetalige vaardigheid \hfill \textbf{Trefwoorden}: Professioneel schrijven, Presenteren}
\cvline{Engels}{Volledige professionele vaardigheid \hfill \textbf{Trefwoorden}: Dagelijkse samenwerking, Rapporteren}
\cvline{Duits}{Beperkte werkvaardigheid \hfill \textbf{Trefwoorden}: Leesvaardigheid, Basisgesprekken}
\section{Vaardigheden}

\cvline{Data-analyse}{Expert \hfill \textbf{Trefwoorden}: SQL, Python, Excel, A/B-testen, Statistiek}
\cvline{Datavisualisatie}{Geavanceerd \hfill \textbf{Trefwoorden}: Power BI, Tableau, Looker Studio, Dashboarding}
\cvline{Data-engineering}{Intermediair \hfill \textbf{Trefwoorden}: dbt, BigQuery, Airflow, Git}
\cvline{Machine Learning}{Intermediair \hfill \textbf{Trefwoorden}: scikit-learn, Regressieanalyse, Beslissingsbomen, Clustering}
\section{Onderscheidingen}

\cventry{dec 2023}
        {Finlytics B.V.}
        {Topscoorder Klanttevredenheid}
        {}
        {}
        {Uitgeroepen tot medewerker met de hoogste stakeholderscore vanwege heldere communicatie en bruikbare analyses.}
\section{Certificaten}

\cventry{jun 2022}
        {Microsoft}
        {Microsoft Certified Power BI Data Analyst Associate}
        {\url{https://learn.microsoft.com/certifications/power-bi-data-analyst-associate}}
        {}
        {}

\cventry{mei 2020}
        {Google via Coursera}
        {Google Data Analytics Certificate}
        {\url{https://www.coursera.org/professional-certificates/google-data-analytics}}
        {}
        {}
\section{Projecten}

\cventry{jan 2023–apr 2023}
        {Machine learning-model om klantverloop in de telecomsector te voorspellen en de belangrijkste risicofactoren te identificeren.}
        {Churnvoorspelling telecomklanten}
        {\url{https://github.com/sannevandenberg/churn-voorspelling}}
        {}
        {\begin{itemize}
\item Bouwde een Random Forest-classificatiemodel in Python met scikit-learn
\item Verbeterde de voorspellingskracht met 18 procent ten opzichte van het bestaande model
\item Leverde een interactief Power BI-dashboard op met de belangrijkste KPI's en risicofactoren
\end{itemize}
\textbf{Trefwoorden}: Python, scikit-learn, Churn, Power BI}
\section{Interesses}

\cvline{Hardlopen}{Halve marathon, Intervaltraining}
\cvline{Koken}{Vegetarisch, Wereldkeuken}
\cvline{Duurzaamheid}{Energietransitie, Circulaire economie}
\section{Vrijwilligerswerk}

\cventry{mrt 2021–Heden}
        {Vrijwillige data-analist}
        {Stichting Data voor Goed}
        {\url{https://www.datavoorgoed.nl}}
        {}
        {\begin{itemize}
\item Help non-profitorganisaties met datagedreven besluitvorming
\item Bouw laagdrempelige dashboards in Power BI
\item Geef workshops over datakwaliteit en privacyvriendelijk datagebruik
\end{itemize}}

\end{document}